% ====================================================================
% CAPÍTULO 3: METODOLOGÍA Y MODELAMIENTO QUÍMICO
% ====================================================================
\chapter{Metodología}
% Nota de avance: Aquí se detalla el modelamiento matemático del 
% potencial eléctrico superficial del nódulo.

% Para cuantificar la disociación espontánea de las moléculas de agua en la interfaz abisal, se emplearon ecuaciones de balance electroquímico adaptadas a altas presiones hidrostáticas. El umbral teórico mínimo de diferencia de potencial necesario para iniciar la hidrólisis galvánica se modela a través de la ecuación \ref{eq:electrolisis}:

\begin{equation}
\label{eq:electrolisis}
2H_2O \rightarrow O_2 + 4H^+ + 4e^-
\end{equation}

% Donde la liberación neta de oxígeno libre ($O_2$) ocurre en la superficie externa del agregado polimetálico toda vez que el gradiente de voltaje superficial medido supere el límite crítico de $1.23 \text{ V}$. Las tasas de flujo másico resultantes de este proceso son posteriormente integradas de forma espacial para calibrar los mapas de distribución geográfica en el fondo marino.